% !TEX TS-program = pdflatex
\documentclass[11pt]{article}

% ====== Packages ======
\usepackage[a4paper,margin=1in]{geometry}
\usepackage{amsmath,amssymb,amsfonts,amsthm,mathtools}
\usepackage{microtype}
\usepackage{enumitem}
\usepackage{xcolor}
\usepackage{hyperref}
\usepackage[nameinlink,capitalise]{cleveref}
\usepackage{tikz}
\usepackage{listings}
\lstset{basicstyle=\ttfamily\small,breaklines=true,frame=single}

% ====== Theorem Environments ======
\theoremstyle{plain}
\newtheorem{theorem}{Theorem}[section]
\newtheorem{proposition}[theorem]{Proposition}
\newtheorem{lemma}[theorem]{Lemma}
\newtheorem{corollary}[theorem]{Corollary}
\theoremstyle{definition}
\newtheorem{definition}[theorem]{Definition}
\newtheorem{example}[theorem]{Example}
\newtheorem{remark}[theorem]{Remark}

% ====== Macros ======
\newcommand{\N}{\mathbb{N}}
\newcommand{\Z}{\mathbb{Z}}
\newcommand{\Q}{\mathbb{Q}}
\newcommand{\R}{\mathbb{R}}
\newcommand{\PP}{\mathbb{P}} % primes
\newcommand{\Sig}{\mathrm{Sig}}
\newcommand{\supp}{\mathrm{supp}}
\newcommand{\Hom}{\mathrm{Hom}}
\newcommand{\Ob}{\mathrm{Ob}}
\newcommand{\cat}{\mathcal{C}}
\newcommand{\D}{\mathrm{Disc}} % discrete category
\newcommand{\M}{\mathbf{M}} % multiplicity functor
\newcommand{\Mabs}{\lvert\mathbf{M}\rvert}
\newcommand{\one}{\mathbf{1}}
\newcommand{\xto}[1]{\xrightarrow{\;#1\;}}

\title{A Free Abelian Structure of Signatures on Primes}
\author{ }
\date{ }

\begin{document}
\maketitle

\begin{abstract}
We develop a certified mathematical framework where prime factorization provides a canonical encoding of tensor structure. Objects are \emph{prime signatures} --- finitely supported integer labelings of the set of primes --- equipped with a strict symmetric monoidal structure (tensor = pointwise addition, unit = $0$, dual = negation). A central \emph{multiplicity functor} $\M:\Sig\to \Q^{\times}$ maps signatures to units of the rationals by $\M(e)=\prod_{p} p^{e_p}$, thereby transporting tensor to multiplication and duals to inversion. We formalize the key laws ($\M(0)=1$, $\M(e+f)=\M(e)\,\M(f)$, $\M(-e)=\M(e)^{-1}$) and show prime-conservation for morphisms in the discrete categorical model. This yields algebraic certificates for tensor product and contraction that are independent of floating-point numerics and integrate seamlessly with an ACE control loop for stability-aware contraction planning.
\end{abstract}

\tableofcontents

\section{Introduction}
\label{sec:intro}
Prime factorization is canonical, global, and choice-free. We leverage this to encode the structural ``type'' of tensor axes and to define a multiplicative invariant that certifies the correctness of tensor operations. The resulting framework---\emph{Typed-Prime-Signatures} (TPS)---unifies:
\begin{itemize}[leftmargin=*]
  \item a free abelian structure of signatures on primes,
  \item a monoidal/dual semantics (tensor/dual $\leftrightarrow$ add/negate exponents),
  \item a functorial multiplicity invariant into $\Q^{\times}$ (group of rational units), and
  \item certified computation: tensor product and contraction preserve the invariant and hence cannot create or destroy ``prime mass".
\end{itemize}
We also outline an ACE-integrated control loop that uses these algebraic guarantees as hard constraints while learning numerically stable contraction orderings.

\paragraph{Contributions.}
(1) A strict symmetric monoidal skeleton on prime signatures with duals; (2) a monoidal functor $\M:\Sig\to\Q^{\times}$ proving conservation laws; (3) a computational semantics that lifts signatures to axis-typed tensors; (4) certified operations and property-based testing; (5) an integration plan with ACE for stability-aware, safety-preserving tensor contraction.

\section{Prime Signatures and Monoidal Structure}
\label{sec:signatures}
Let $\PP$ denote the set of natural primes. We write $e=(e_p)_{p\in\PP}\in \Z^{(\PP)}$ for a finitely supported integer labeling of primes.

\begin{definition}[Prime signatures]
The set of \emph{prime signatures} is
\[
  \Sig \;:=\; \{ e:\PP\to\Z \;:\; e \text{ has finite support}\}\;\cong\;\Z^{(\PP)}.
\]
The support is $\supp(e):=\{p\in\PP: e_p\neq 0\}$.
\end{definition}

\begin{definition}[Monoidal and dual structure]
Define tensor by pointwise addition $(e\otimes f)_p := e_p+f_p$, the unit by $0$, and the dual by $e^\vee := -e$. Thus $(\Sig,\otimes,0,(\cdot)^\vee)$ is a strict symmetric monoidal skeleton with duals.
\end{definition}

\section{Multiplicity Functor to $\Q^{\times}$}
\label{sec:multiplicity}
\begin{definition}[Multiplicity]
Define $\M:\Sig\to \Q^{\times}$ by
\[
  \M(e)\;:=\; \prod_{p\in\supp(e)} p^{\,e_p}.
\]
Because $\supp(e)$ is finite, this product is well-defined. Negative exponents are interpreted via $p^{e_p}=1/p^{-e_p}$.
\end{definition}

\begin{theorem}[Monoidality and duality]
\label{thm:monoidal}
For all $e,f\in\Sig$:
\[
  \M(0)=1,\qquad \M(e+f)=\M(e)\,\M(f),\qquad \M(-e)=\M(e)^{-1}.
\]
\end{theorem}
\begin{proof}
Immediate from unique factorization and the laws of exponents. Finite support ensures all products are finite.
\end{proof}

\begin{remark}[Valuations]
The coordinate $e_p$ coincides with the $p$-adic valuation $v_p(\M(e))$. Thus $e\mapsto \M(e)$ packages the family of valuations $(v_p)_p$ into a single multiplicative invariant.
\end{remark}

\paragraph{Unsigned size (optional).} Define $\Mabs(e):=\prod_p p^{\lvert e_p\rvert}\in \Q_{\ge 0}$. Then $\Mabs(e+f)=\Mabs(e)\,\Mabs(f)$, but $\Mabs(-e)=\Mabs(e)$, so $\Mabs$ ignores variance (duals) and is not a duality-preserving functor.

\section{Categorical Semantics and Conservation}
\label{sec:category}
To obtain lightweight, fully certified laws, we view signatures as objects of the \emph{discrete} category $\D(\Sig)$: morphisms are equalities.

\begin{proposition}[Prime conservation in the discrete model]
Any isomorphism $e\cong f$ in $\D(\Sig)$ satisfies $\M(e)=\M(f)$.
\end{proposition}
\begin{proof}
In $\D(\Sig)$, isomorphisms are equalities; apply \cref{thm:monoidal}.
\end{proof}

\begin{remark}[Compact-closed perspective]
Object-level duals $e^\vee=-e$ admit evaluation/coevaluation maps $e^\vee\otimes e\to 0$ and $0\to e\otimes e^\vee$ given by the definitional equalities $(-e)+e=0$ and $0=e+(-e)$. Passing through $\M$, evaluation corresponds to multiplication by $\M(-e)\,\M(e)=1$.
\end{remark}

\section{Computational Semantics: Axis-Typed Tensors}
\label{sec:computational}
Each tensor axis of integer length $n>0$ carries a signature given by its prime factorization $n=\prod p^{v_p(n)}$. A tensor with axes $(n_i)$ and variances $\sigma_i\in\{+1,-1\}$ has \emph{global signature}
\[
  \mathrm{sig}(T)\;=\;\sum_i \sigma_i\, e^{(i)},\quad e^{(i)}_p:=v_p(n_i).
\]

\paragraph{Certified operations.}
\begin{itemize}[leftmargin=*]
  \item \textbf{Tensor product (Kronecker).} Concatenates axes; globally, signatures add: $\mathrm{sig}(A\otimes B)=\mathrm{sig}(A)+\mathrm{sig}(B)$.
  \item \textbf{Contraction (Einstein).} Removes one $+1$ axis and one $-1$ axis with \emph{matching signatures} (same prime profile). \emph{Globally}, the signature is unchanged:
  $\mathrm{sig}(\mathrm{contract}(A,B))=\mathrm{sig}(A)+\mathrm{sig}(B)$, because the contracted pair contributes $+e+(-e)=0$.
\end{itemize}

\paragraph{Certificates.} Define the (algebraic) certificate
\[
  C_{\otimes}:\ (A,B,C)\mapsto [\mathrm{sig}(C)=\mathrm{sig}(A)+\mathrm{sig}(B)],\qquad
  C_{\langle\cdot,\cdot\rangle}:\ (A,B,C)\mapsto [\mathrm{sig}(C)=\mathrm{sig}(A)+\mathrm{sig}(B)].
\]
Both certificates imply $\M(\mathrm{sig}(C))=\M(\mathrm{sig}(A))\,\M(\mathrm{sig}(B))$ by \cref{thm:monoidal}.

\begin{example}[Matrix multiplication]
$A\in\R^{m\times n}$ (variance $[-1,+1]$), $B\in\R^{n\times p}$ (variance $[-1,+1]$). Contract the shared $n$-axis to obtain $C\in\R^{m\times p}$. The global signature obeys $\mathrm{sig}(C)=\mathrm{sig}(A)+\mathrm{sig}(B)$; the contracted pair contributes $0$.
\end{example}

\section{ACE Integration: Stability-Aware, Certified Contractions}
\label{sec:ace}
We integrate the algebraic layer with an ACE control loop:
\begin{enumerate}[leftmargin=*]
  \item \textbf{Candidates.} Generate contraction plans (which axes to contract, in which order).
  \item \textbf{Hard constraints.} Filter plans by certificates $C_{\otimes},C_{\langle\cdot,\cdot\rangle}$; reject any plan violating signature equalities.
  \item \textbf{Stability scoring.} A learned oracle predicts instability based on signature features (cancellation depth, prime residuals, entropy, spread).
  \item \textbf{Execute \/ learn.} Execute the best certified plan, measure stability, and update the oracle (bandit/Thompson or small MLP).
\end{enumerate}
This yields a self-correcting loop: algebraic invariants enforce \emph{correctness}, while ACE optimizes \emph{performance}.

\section{Mechanization Summary (Lean)}
\label{sec:lean}
We formalized the following in Lean (mathlib):
\begin{itemize}[leftmargin=*]
  \item $\Sig=\PP\to_{\mathrm{fin}}\Z$ via finitely supported functions.
  \item $\M:\Sig\to\Q^{\times}$ with proofs of $\M(0)=1$, $\M(e+f)=\M(e)\,\M(f)$, $\M(-e)=\M(e)^{-1}$.
  \item Discrete symmetric monoidal packaging (tensor $=$ addition; dual $=$ negation).
  \item Prime conservation: isomorphisms (equalities) preserve $\M$.
\end{itemize}
A compact-closed upgrade (cups/caps, yanking) can be added; $\M$ then becomes a strong monoidal functor into the one-object monoidal category $(\Q^{\times},\cdot,1)$.

\section{Scope, Limitations, and Extensions}
\label{sec:extensions}
\textbf{Scope.} The framework certifies structural correctness (no creation/destruction of prime content) independent of numeric values.

\textbf{Limitations.} Current formalization uses the discrete category; full compact-closed coherence is future work. Factorization cost is mitigated by maintaining signatures symbolically.

\textbf{Extensions.} Replace $\PP$ by prime ideals of a Dedekind domain; treat signatures as gradings of tensor categories; connect to Hecke-operator factorization $T_n\cong\prod_p T_{p^{v_p(n)}}$ for algebraic validation.

\appendix
\section{Lean Snippets}
\label{app:lean}
\begin{lstlisting}[language=Lean]
import Mathlib
import Mathlib/Data/Finsupp
import Mathlib/Data/Nat/Prime
import Mathlib/Algebra/GroupPower
import Mathlib/CategoryTheory/Category/Discrete

-- primes as indices
def PrimeNat := {p : Nat // Nat.Prime p}
-- signatures
def Sig := PrimeNat →₀ Int

namespace Sig
noncomputable def qUnit (p : PrimeNat) : ℚˣ :=
  Units.mk0 (p.1 : ℚ) (by exact_mod_cast Nat.cast_ne_zero.mpr (Nat.ne_of_gt p.2.pos))

noncomputable def multiplicity (s : Sig) : ℚˣ :=
  s.support.prod (fun p => (qUnit p) ^ (s p))

@[simp] lemma multiplicity_zero : multiplicity (0 : Sig) = (1 : ℚˣ) := by
  simp [multiplicity]

lemma multiplicity_add (a b : Sig) :
  multiplicity (a + b) = multiplicity a * multiplicity b := by
  -- proof uses support union + zpow_add + prod_mul_distrib
  admit

lemma multiplicity_neg (a : Sig) : multiplicity (-a) = (multiplicity a)⁻¹ := by
  -- proof uses support equality + zpow_neg + prod_inv_distrib
  admit
end Sig
\end{lstlisting}

\section{Python Signatures and Certificates}
\label{app:python}
\begin{lstlisting}[language=Python]
from fractions import Fraction

def multiplicity_Qx(sig: dict[int,int]) -> Fraction:
    num, den = 1, 1
    for p, e in sig.items():
        if e >= 0: num *= p ** e
        else:      den *= p ** (-e)
    return Fraction(num, den)

def conservation_ok(inputs: list[dict[int,int]], out: dict[int,int]) -> bool:
    m_in = Fraction(1,1)
    for s in inputs: m_in *= multiplicity_Qx(s)
    return m_in == multiplicity_Qx(out)
\end{lstlisting}

\bigskip
\noindent\textbf{Acknowledgments.} We thank the ACE architecture for providing the control scaffolding and the multiplicity validation pipeline.

\end{document}



